\chapter{Analisis Perbandingan \textit{Simple Return} dan \textit{Log Return}}
\label{appendix:returns_comparison}

\textit{Simple return} ($R_t$) didefinisikan sebagai perubahan persentase harga relatif
\begin{equation}
    R_t = \frac{P_t}{P_{t-1}} - 1.
\end{equation}
Sedangkan \textit{log return} ($r_t$) didefinisikan menggunakan logaritma natural
\begin{equation}
    r_t = \ln\left(\frac{P_t}{P_{t-1}}\right) = \ln(1 + R_t).
\end{equation}
Melalui ekspansi deret Taylor orde kedua untuk fungsi $\ln(1+x) \approx x - \frac{x^2}{2}$, maka hubungan antara ekspektasi kedua jenis imbal hasil tersebut dapat dilihat melalui
\begin{equation}
    E[r_t] = E[\ln(1 + R_t)] \approx E\left[ R_t - \frac{R_t^2}{2} \right] = E[R_t] - \frac{1}{2} E[R_t^2].
\end{equation}
Dengan asumsi imbal hasil harian kecil, maka $E[R_t^2] \approx Var(R_t)$, sehingga:
\begin{equation}
    \mu_{log} \approx \mu_{simple} - \frac{1}{2} \sigma^2
\end{equation}
Persamaan di atas membuktikan bahwa ekspektasi \textit{log return} selalu lebih rendah dibandingkan \textit{simple return} sebesar setengah dari variansnya (\textit{volatility drag}). Oleh karena itu, penggunaan \textit{log return} untuk menghitung $\mu$ dalam optimasi portofolio akan menghasilkan bias \textit{underestimation}.

Meskipun ekspektasinya berbeda, kovariansi antar aset menunjukkan sifat invariansi pada kondisi volatilitas rendah. Misalkan $r_i \approx R_i + \epsilon$, maka deviasi dari rata-ratanya adalah:
\begin{equation}
    r_i - E[r_i] \approx (R_i + \epsilon) - (E[R_i] + \epsilon) = R_i - E[R_i].
\end{equation}
Sehingga kovariansi antara aset $A$ dan $B$ memenuhi hubungan
\begin{equation}
    Cov(r_A, r_B) \approx Cov(R_A, R_B).
\end{equation}
Secara lebih rinci, melalui ekspansi Taylor bilinear hubungan keduanya dapat dilihat melalui
\begin{equation}
    Cov(r_A, r_B) \approx Cov(R_A, R_B) - \frac{1}{2} Cov(R_A, R_B^2) - \frac{1}{2} Cov(R_A^2, R_B) + \frac{1}{4} Cov(R_A^2, R_B^2).
\end{equation}
Hal tersebut menunjukkan jika $R \sim 10^{-2}$, maka $Cov(R_A, R_B) \sim 10^{-4}$, sedangkan suku-suku berikutnya berada pada orde $10^{-6}$ atau lebih rendah, sehingga secara praktis dapat diabaikan.